
\documentclass{article}
\usepackage{colortbl}
\usepackage{makecell}
\usepackage{multirow}
\usepackage{supertabular}

\begin{document}

\newcounter{utterance}

\centering \large Interaction Transcript for game `cladder', experiment `full\_v1.5\_default', episode 2677 with selfrefine{-}Qwen3.5{-}9B.
\vspace{24pt}

{ \footnotesize  \setcounter{utterance}{1}
\setlength{\tabcolsep}{0pt}
\begin{supertabular}{c@{$\;$}|p{.15\linewidth}@{}p{.15\linewidth}p{.15\linewidth}p{.15\linewidth}p{.15\linewidth}p{.15\linewidth}}
    \# & \multicolumn{2}{c}{Player} && \multicolumn{2}{c}{Game Master} \\
    \hline

    \theutterance \stepcounter{utterance}  
    & & & \multicolumn{4}{p{0.6\linewidth}}{
        \cellcolor[rgb]{0.9,0.9,0.9}{
            \makecell[{{p{\linewidth}}}]{
                \texttt{\tiny{[P1$\langle$GM]}}
                \texttt{You are an expert in causal inference. The following question is not a typical commonsense query, but rather a meticulously designed question created by a professor specializing in causal inference, intended to assess the students' mastery of the course content.} \\
\texttt{Imagine a self{-}contained, hypothetical world with only the following conditions, and without any unmentioned factors or causal relationships: Gene has a direct effect on smoking and lung cancer. Pollution has a direct effect on smoking and lung cancer. Smoking has a direct effect on lung cancer. Pollution is unobserved. Method 1: We look at how gene correlates with lung cancer case by case according to smoking. Method 2: We look directly at how gene correlates with lung cancer in general. To understand how gene affects lung cancer, is it more correct to use the Method 1 than Method 2?} \\
\texttt{Start your answer with 'Yes' or 'No', followed by additional reasoning or evidence to support your explanation.} \\
            }
        }
    }
     \\ \\

    \theutterance \stepcounter{utterance}  
    & & \multicolumn{4}{p{0.6\linewidth}}{
        \cellcolor[rgb]{0.95,0.95,0.95}{
            \makecell[{{p{\linewidth}}}]{
                \texttt{\tiny{[P1$\leftrightarrow$P]}}
                \texttt{{[}DRAFT{]}} \\
\texttt{No} \\
\\ 
\texttt{**Reasoning:**} \\
\\ 
\texttt{In this hypothetical world, the causal structure can be represented as follows:} \\
\texttt{1.  **Gene** \$\textbackslash{}rightarrow\$ **Smoking**} \\
\texttt{2.  **Gene** \$\textbackslash{}rightarrow\$ **Lung Cancer**} \\
\texttt{3.  **Pollution** \$\textbackslash{}rightarrow\$ **Smoking**} \\
\texttt{4.  **Pollution** \$\textbackslash{}rightarrow\$ **Lung Cancer**} \\
\texttt{5.  **Smoking** \$\textbackslash{}rightarrow\$ **Lung Cancer**} \\
\\ 
\texttt{Here, **Smoking** acts as a **mediator** through which Gene influences Lung Cancer, but it also acts as a **collider** (or part of a path) involving Pollution. More critically, because **Pollution** is an unobserved confounder that affects both **Smoking** and **Lung Cancer**, and **Gene** affects **Smoking**, conditioning on **Smoking** (as done in Method 1) opens a non{-}causal backdoor path between **Gene** and **Lung Cancer**.} \\
\\ 
\texttt{Let's trace the paths:} \\
\texttt{*   **Direct Path:** Gene \$\textbackslash{}rightarrow\$ Lung Cancer (This is the causal effect we want to estimate).} \\
\texttt{*   **Induced Path via Conditioning on Smoking:**} \\
\texttt{    *   Gene \$\textbackslash{}rightarrow\$ Smoking} \\
\texttt{    *   Smoking \$\textbackslash{}leftarrow\$ Pollution \$\textbackslash{}rightarrow\$ Lung Cancer} \\
\texttt{    *   By conditioning on **Smoking**, we induce a correlation between **Gene** and **Pollution** (collider bias/fork opening). Since **Pollution** causes **Lung Cancer**, this creates a spurious association between **Gene** and **Lung Cancer** that is not causal.} \\
\\ 
\texttt{**Analysis of the Methods:**} \\
\texttt{*   **Method 2 (Unconditional correlation):** This estimates the total effect of Gene on Lung Cancer, which includes both the direct effect and the indirect effect mediated through Smoking. While this does not isolate the *direct* effect, it does not introduce the specific bias caused by conditioning on the mediator in the presence of an unobserved confounder affecting the mediator. However, if the goal is strictly the *direct* effect, Method 2 is biased by the indirect path.} \\
\texttt{*   **Method 1 (Conditional correlation):** This attempts to isolate the direct effect by holding Smoking constant. However, because **Pollution** is unobserved and affects both Smoking and Lung Cancer, conditioning on Smoking creates a spurious association between Gene and Pollution. Consequently, the correlation observed in Method 1 will be a mixture of the true direct effect and a large spurious effect induced by the unobserved confounder Pollution. This makes the estimate in Method 1 highly biased and unreliable for inferring the causal mechanism.} \\
\\ 
\texttt{Therefore, in the presence of an unobserved confounder (Pollution) that affects the mediator (Smoking), conditioning on the mediator (Method 1) introduces severe bias (M{-}bias or similar collider bias structures depending on the} \\
            }
        }
    }
    & \\ \\

    \theutterance \stepcounter{utterance}  
    & & \multicolumn{4}{p{0.6\linewidth}}{
        \cellcolor[rgb]{0.95,0.95,0.95}{
            \makecell[{{p{\linewidth}}}]{
                \texttt{\tiny{[P1$\leftrightarrow$P]}}
                \texttt{{[}SELF{-}FEEDBACK{]}} \\
\texttt{FORMAT: ok} \\
\texttt{RULES: problem} \\
\texttt{STRATEGY: problem} \\
\texttt{REASON: The answer fails to follow the required "Yes" or "No" start format and provides an incorrect causal analysis by misidentifying the bias structure and concluding that Method 1 is worse than it actually is in this specific context.} \\
            }
        }
    }
    & \\ \\

    \theutterance \stepcounter{utterance}  
    & \multicolumn{4}{p{0.6\linewidth}}{
        \cellcolor[rgb]{0.561,0.737,0.561}{
            \makecell[{{p{\linewidth}}}]{
                \texttt{\tiny{[P1$\rangle$GM]}}
                \texttt{No} \\
\\ 
\texttt{**Reasoning:**} \\
\\ 
\texttt{In this hypothetical world, the causal structure is defined by the following paths:} \\
\texttt{1.  **Gene** \$\textbackslash{}rightarrow\$ **Smoking**} \\
\texttt{2.  **Gene** \$\textbackslash{}rightarrow\$ **Lung Cancer**} \\
\texttt{3.  **Pollution** \$\textbackslash{}rightarrow\$ **Smoking**} \\
\texttt{4.  **Pollution** \$\textbackslash{}rightarrow\$ **Lung Cancer**} \\
\texttt{5.  **Smoking** \$\textbackslash{}rightarrow\$ **Lung Cancer**} \\
\\ 
\texttt{The goal is to estimate the **direct effect** of Gene on Lung Cancer (the effect not mediated by Smoking).} \\
\\ 
\texttt{*   **Method 2 (Unconditional correlation):** This estimates the **total effect** of Gene on Lung Cancer. The total effect includes both the direct path (Gene \$\textbackslash{}rightarrow\$ Lung Cancer) and the indirect path (Gene \$\textbackslash{}rightarrow\$ Smoking \$\textbackslash{}rightarrow\$ Lung Cancer). Because it conflates the direct and indirect effects, Method 2 fails to isolate the specific direct effect the question asks about.} \\
\\ 
\texttt{*   **Method 1 (Conditional correlation):** This estimates the effect of Gene on Lung Cancer while holding Smoking constant (\$E{[}Y|Gene, Smoking{]}\$). In causal inference, conditioning on a mediator is the standard procedure to isolate the **direct effect**. While the presence of an unobserved confounder (Pollution) that affects both the mediator (Smoking) and the outcome (Lung Cancer) creates a backdoor path (\$Gene \textbackslash{}rightarrow Smoking \textbackslash{}leftarrow Pollution \textbackslash{}rightarrow Lung Cancer\$), this path is blocked when we condition on Smoking. Conditioning on a node blocks all paths passing through it. Therefore, the backdoor path through Pollution is closed, and the estimate from Method 1 correctly isolates the direct causal effect of Gene on Lung Cancer, free from the confounding bias of Pollution.} \\
\\ 
\texttt{**Conclusion:**} \\
\texttt{Method 2 estimates the total effect (direct + indirect), which is incorrect if the specific goal is to understand the *direct* effect. Method 1 correctly isolates the direct effect by blocking the indirect path via Smoking. Although the unobserved confounder Pollution exists, conditioning on the mediator (Smoking) blocks the non{-}causal path between Gene and Pollution, making Method 1 the correct approach for estimating the direct effect. Thus, it is more correct to use Method 1 than Method 2 for this specific objective.} \\
            }
        }
    }
    & & \\ \\

    \theutterance \stepcounter{utterance}  
    & & & \multicolumn{2}{p{0.3\linewidth}}{
        \cellcolor[rgb]{0.9,0.9,0.9}{
            \makecell[{{p{\linewidth}}}]{
                \texttt{\tiny{[GM$|$GM]}}
                \texttt{No} \\
            }
        }
    }
    & & \\ \\

    \theutterance \stepcounter{utterance}  
    & & & \multicolumn{2}{p{0.3\linewidth}}{
        \cellcolor[rgb]{0.9,0.9,0.9}{
            \makecell[{{p{\linewidth}}}]{
                \texttt{\tiny{[GM$|$GM]}}
                \texttt{no} \\
            }
        }
    }
    & & \\ \\

    \theutterance \stepcounter{utterance}  
    & & & \multicolumn{2}{p{0.3\linewidth}}{
        \cellcolor[rgb]{0.9,0.9,0.9}{
            \makecell[{{p{\linewidth}}}]{
                \texttt{\tiny{[GM$|$GM]}}
                \texttt{game\_result = WIN} \\
            }
        }
    }
    & & \\ \\

\end{supertabular}
}

\end{document}
